\writeSectionFrame{\presentationTitle}{Das Problem}
\begin{frame}{Zustandsautomaten}
	\begin{center}
 		\includegraphics[scale=0.3]{\BILDERPFAD/problem1.png}
 	\end{center}
\end{frame}

\note{
	Das Zustand-Muster ist eng verwandt mit einem Zustandsautomaten. Es gibt eine endliche
	Menge an Zuständen, in denen das Programm sein kann. In einem Zustand verhält sich das
	Programm jeweils unterschiedlich und es kann von einem Zustand in einen bestimmten
	anderen Zustand übergehen.  
}

\begin{frame}{Änderung von Zuständen}
	\begin{alertblock}{Das ist doch was normales}
 		if (einAttribut > einVergleichswert)
 	\end{alertblock}
\end{frame}

\note{
	Aber was ist jetzt besonders daran? Der Zustand von Objekten ändert sich doch dauernd. Und dabei
	ändert sich auch das Verhalten eines Objekts. Normalerweise wird dazu irgendeine if-Anweisung
	ausgewertet. Aber manchmal ändert sich das Verhalten eines Objekts so grundlegend, dass man
	einen Teil des Objekts austauschen möchte. 
}

\begin{frame}{Thread-Verwaltung}
	\begin{center}
 		\includegraphics[scale=0.27]{\BILDERPFAD/problem2.png}
 	\end{center}
\end{frame}

\note{
	Ein gutes Beispiel sind Threads. Ein Thread durchläuft verschiedene Zustände und es leuchtet 
	sofort ein, dass ein Thread, der wartet, sich ganz anders verhält als ein laufender Thread. Man könnte dies
	mit Fallunterscheidungen und Zustandsvariablen kontrollieren. Das würde die Lesbarkeit und vor allem die	Änderbarkeit nicht erhöhen. 
}